\chapter{乘积测度}

\begin{solution}{product-sigma-coordinate}
对 \(A\in\mathcal A\)，
\(\pi_X^{-1}(A)=A\times Y\in\mathcal A\otimes\mathcal B\)，故
\(\pi_X\) 可测；另一坐标同理。对目标空间中的可测矩形
\(A'\times B'\)，乘积映射的原像为
\(f^{-1}(A')\times g^{-1}(B')\)，属于 \(\mathcal A\otimes\mathcal B\)。
原像保持补和可数并，而这些矩形生成目标的乘积 \(\sigma\)-代数，
所以乘积映射可测。
\end{solution}

\begin{solution}{borel-product-euclidean}
坐标投影连续，因此对任意 Borel 集 \(A\subseteq\R^m\)、
\(B\subseteq\R^n\)，矩形
\(A\times B=\pi_1^{-1}(A)\cap\pi_2^{-1}(B)\) 是 Borel 集。故
\(\mathcal B(\mathbb R^m)\otimes\mathcal B(\mathbb R^n)
\subset\mathcal B(\mathbb R^{m+n})\)。反过来，有理中心、有理半径的开球
构成可数基，每个开集是可数个基元素之并，而每个开球是可数个有理端点开
矩形之并。因此每个开集属于乘积 \(\sigma\)-代数。
\end{solution}

\begin{solution}{product-rectangle-refinement}
收集所有第一坐标集合 \(A_k\)，由它们及补集形成有限个两两不交原子
\(C_i\)；对所有 \(B_k\) 同样形成 \(D_j\)。每个原矩形都是若干
\(C_i\times D_j\) 的并，所以矩形有限并可细分为两两不交矩形。若同一
集合有两种表示，取两种表示的共同原子网格；在每个原子上，两种表示的指示
值相同，有限可加性使两边的
\(\sum\mu(C_i)\nu(D_j)\) 相同。
\end{solution}

\begin{solution}{section-algebra}
逐元素验证：
\[
 y\in(E^c)_x\iff(x,y)\notin E\iff y\notin E_x,
\]
以及
\[
 y\in(\cup_nE_n)_x\iff\exists n,\ (x,y)\in E_n
\iff y\in\cup_n(E_n)_x.
\]
所以“所有截面可测”的集合族对补和可数并封闭，是一个包含矩形的
\(\sigma\)-代数。
\end{solution}

\begin{solution}{cavalieri-graph}
映射 \(F(x,y)=y-f(x)\) 可测，故
\(\Gamma_f=F^{-1}(\{0\})\) 是 Borel 集。每个竖截面
\((\Gamma_f)_x=\{f(x)\}\) 的一维测度为零。Cavalieri 公式给出
\[
 m_2(\Gamma_f)=\int_\mathbb R m_1(\{f(x)\})\,dx=0.
\]
\end{solution}

\begin{solution}{tonelli-layer-cake}
令 \(E=\{(x,t):0<t<f(x)\}\)。函数
\((x,t)\mapsto f(x)-t\) 可测，故 \(E\) 可测。对固定 \(x\)，
\(E_x=(0,f(x))\)，长度为 \(f(x)\)；对固定 \(t\)，
\(E^t=\{f>t\}\)。Tonelli 定理应用于 \(\mathbf1_E\) 得
\[
 \int_Xf\,d\mu
 =\int_{X\times(0,\infty)}\mathbf1_E
 =\int_0^\infty\mu(\{f>t\})\,dt.
\]
\end{solution}

\begin{solution}{tonelli-region-integral}
固定 \(x\in(0,1)\)，竖截面为 \((0,1-x)\)，故
\[
 m_2(E)=\int_0^1(1-x)\,dx=\frac12.
\]
固定 \(y\in(0,1)\) 时横截面为 \((0,1-y)\)，交换次序得到同样结果。
\end{solution}

\begin{solution}{fubini-section-null}
由集合形式 Tonelli，
\[
 0=(\mu\times\nu)(E)=\int_X\nu(E_x)\,d\mu.
\]
非负函数零积分判据给出 \(\nu(E_x)=0\) 对几乎所有 \(x\) 成立。若每个
截面都零测且 \(E\) 可测，则同一公式确实推出 \(E\) 零测；若删去
\(E\) 的乘积可测性，截面信息不足以调用该公式。
\end{solution}

\begin{solution}{fubini-integrability-test}
令 \(f(0,0)=0\)，所得函数 Borel 可测。在区域
\(0<x<1/2\)、\(0<y<x/2\) 上，
\[
 f(x,y)\geq\frac{3x^2/4}{(5x^2/4)^2}
 =\frac{12}{25x^2}.
\]
对 \(|f|\) 用 Tonelli 定理，得到
\[
 \int_{(-1,1)^2}|f|\geq
 \int_0^{1/2}\int_0^{x/2}\frac{12}{25x^2}\,dy\,dx
 =\frac6{25}\int_0^{1/2}\frac{dx}{x}=\infty.
\]
这里一元积分的发散可先在 \([\delta,1/2]\) 上计算，再用单调收敛。
固定 \(x\ne0\)，
\[
 \int_{-1}^{1}\frac{x^2-y^2}{(x^2+y^2)^2}\,dy
 =\left[\frac{y}{x^2+y^2}\right]_{-1}^{1}
 =\frac{2}{x^2+1},
\]
在 \(x=0\) 的截面上，函数为 \(-1/y^2\)（\(y\ne0\)），内层反常积分
发散至 \(-\infty\)。因此按题意，先对非零外层参数计算，再取
\[
 \lim_{\delta\downarrow0}\left(\int_{-1}^{-\delta}
 +\int_\delta^1\right)\frac2{1+x^2}\,dx=\pi.
\]
两侧外层积分各自收敛，此处不依赖主值抵消。交换次序后，\(y\ne0\) 时
内层积分为 \(-2/(1+y^2)\)，而 \(y=0\) 的截面为 \(1/x^2\)，发散至
\(+\infty\)。相同的外层极限为 \(-\pi\)。把两个异常参数处的内层值
分别补为零，也给出这两个外层Lebesgue积分；这不使原函数绝对可积，
因而不能调用 Fubini 定理换序。
\end{solution}

\begin{solution}{double-series-counterexample}
固定 \(m\)，只有 \(a_{m,m}=1,a_{m,m+1}=-1\)，行和为零，故先按
\(n\) 求和再按 \(m\) 求和得零。固定 \(n=0\)，列中只有
\(a_{0,0}=1\)；固定 \(n\ge1\)，有
\(a_{n,n}=1,a_{n-1,n}=-1\)，列和为零。因此反向迭代和为 \(1\)。
每行绝对值和为 \(2\)，故绝对值二重和为无穷。
\end{solution}

\begin{solution}{completed-product-exceptional-section}
集合 \(Z=\{0\}\times[0,1]\) 是二维 Borel 零测集，其任意子集均
Lebesgue 可测，所以 \(E=\{0\}\times V\) 可测且测度为零。
若 \(E\in\mathcal L(\R)\otimes\mathcal L(\R)\)，截面可测性命题将推出
\(E_0=V\in\mathcal L(\R)\)，与选择矛盾。
另一方面 \(Z\) 本身属于这个乘积 \(\sigma\)-代数且乘积测度为零。
其子集 \(E\) 不在其中，故两个完备 Lebesgue 测度的乘积测度不完备。
完备化后 \(E\) 已经允许作为积分对象，却有一个不可测截面。
所以完备化版本只能保证几乎每个截面可测。
\end{solution}

\begin{solution}{convolution-measurability}
先设 \(f,g\) Borel 可测。映射 \((x,y)\mapsto x-y\) 连续，所以
\((x,y)\mapsto f(x-y)g(y)\) 非负 Borel 可测，约定 \(0\cdot\infty=0\)。
Tonelli 定理说明对 \(y\) 的积分是 \(x\) 的可测函数。
一般 Lebesgue 可测情形，取非负 Borel 代表 \(f_0,g_0\)，使
\(f=f_0\)、\(g=g_0\) 分别在零测集 \(N_f,N_g\) 以外成立。
对每个固定 \(x\)，两种被积函数只可能在
\((x-N_f)\cup N_g\) 上不同，该集合仍为零测集。
因此 \(f*g=f_0*g_0\) 对每个 \(x\) 成立，卷积可测。
对 Borel 代表再次使用 Tonelli，并利用平移不变性。令
\(A=\int f\)、\(B=\int g\)。对每个固定的 \(y\)，有
\(\int_x f(x-y)g(y)\,dx=A g(y)\)：有限常数时由齐次性；
若 \(g(y)=\infty\)，用常数截断和单调收敛得到同一结论。
于是
\[
 \int_x(f*g)(x)\,dx
 =\int_y g(y)\int_x f(x-y)\,dx\,dy
 =\left(\int f\right)\left(\int g\right).
\]
最后一步的无穷值约定也可直接核验。\(A=0\) 时 \(Ag=0\)；
\(0<A<\infty\) 时用有限常数齐次性；\(A=\infty\) 时，若 \(B=0\)，
则 \(g=0\) 几乎处处，积分为零；若 \(B>0\)，则 \(\{g>0\}\) 有正测度，
\(Ag\) 在该集合上为无穷，积分为无穷。各情形都符合 \(0\cdot\infty=0\)。
\end{solution}

\begin{solution}{associativity-product-measure}
两种测度都定义在
\(\mathcal A\otimes\mathcal B\otimes\mathcal C\) 上，并在矩形
\(A\times B\times C\) 上取值
\(\mu(A)\nu(B)\lambda(C)\)。三测度的 \(\sigma\)-有限性使矩形预测度
\(\sigma\)-有限，延拓唯一性给出两种括号方式相同。对非负可测函数先由
指示函数、简单函数和单调收敛依次得到三重 Tonelli；对绝对可积函数应用于
绝对值，再拆正负部，得到所有积分次序相同的 Fubini 公式。
\end{solution}
